% !PRINTED 514
\weberpage{514}
%<<CONT
n\"amlich durch Adjunction von algebraischen Gr\"ossen m\"oglichst
einfacher Natur eine Zerf\"allung der \gesperrt{Galois}'schen Resolvente,
also eine Erniedrigung des Grades der Gruppe herbeizuf\"uhren.

Wir stellen jetzt die Frage so. In dem urspr\"unglichen
K\"orper $\Omega$ ist die \gesperrt{Galois}'sche Resolvente $g(t)$ irreducibel. Der
K\"orper $\Omega$ soll zu einem anderen K\"orper $\Omega'$ so erweitert werden,
dass $g(t)$ reducibel wird. $\Omega'$ soll dabei ein algebraischer K\"orper
\"uber $\Omega$ sein, und es muss also nach \S.\,143 eine algebraische
Gr\"osse $\varepsilon$ geben, so dass $\Omega' = \Omega(\varepsilon)$ wird. Diese Gr\"osse $\varepsilon$ wird
Wurzel einer gewissen irreducibeln Gleichung in $\Omega$ sein, die
wir mit
\weq{1}{\chi(\varepsilon) = 0}
bezeichnen wollen.

Nehmen wir an, in dem K\"orper $\Omega(\varepsilon)$ sondere sich von $g(t)$
der irreducible Factor
\weq{2}{g_1(t) = g_1(t, \varepsilon)}
ab, der die Wurzel $t = \varrho$ hat. Den Grad von $g_1(t, \varepsilon)$ wollen wir
mit $q$ bezeichnen und die Wurzeln mit
\weq{3}{\varrho, \varrho_1, \varrho_2 \ldots \varrho_{q-1};}
so dass
\weq{4}{g_1(t, \varepsilon) = (t - \varrho)(t - \varrho_1) \ldots (t - \varrho_{q-1}).}
Zun\"achst ist nun nachzuweisen, dass die in der Gruppe $P$
enthaltenen Substitutionen
\[ (\varrho, \varrho), (\varrho, \varrho_1), (\varrho, \varrho_2) \ldots (\varrho, \varrho_{q-1}) \]
eine Gruppe $Q$ bilden. Um dies zu zeigen, setzen wir wie
fr\"uher (\S.\,146)
\[ \varrho_i = \Theta_i(\varrho) \]
und bedenken, dass dann die Gleichung
\[ g_1[\Theta_i(t), \varepsilon] = 0 \]
eine Wurzel, n\"amlich $t = \varrho$ mit $g_1(t)$ gemein hat und mithin,
da $g_1(t)$ irreducibel ist, durch $g_1(t)$ theilbar ist. Daraus folgt,
dass, wenn $\varrho_1$ und $\varrho_2$ irgend zwei der Wurzeln (3) sind, auch
$\Theta_i(\varrho_2) = \varrho_3$ unter diesen Wurzeln enthalten ist. Es ist aber die
Substitution $(\varrho, \varrho_3)$ aus $(\varrho, \varrho_1)$ und $(\varrho, \varrho_2)$ zusammengesetzt, und
damit ist die Gruppeneigenschaft von $Q$ nachgewiesen. Wir bezeichnen den Index von $Q$ mit $j$ und setzen
\weq{5}{p = jq.}
